\section{Experiments}\label{sec:experiments}
\subsection{Datasets and Few-Shot Splits}
We evaluate MVTec AD~\cite{mvtec2019} and VisA~\cite{visa2022}. The strict external-transfer audit additionally incorporates MPDD\@. MVTec AD comprises industrial object and texture categories with pixel-level defect masks. VisA contains a broader range of industrial categories and serves as the primary benchmark for full-scale corruption and calibration shift. MPDD provides six metal-part categories for an external stress test that transfers from MVTec to MPDD. Within-MPDD certification is not pursued because the frozen nested implementation requires at least six eligible source categories, whereas excluding one MPDD target leaves five. For each category and seed, one random permutation of the path-sorted normal training images defines nested support sets, so the $k$-shot set is contained in every larger declared set. The CRESS source archive is separate: it consists of held-out known-normal evaluation images from non-target source categories (the \texttt{test/good} partition for MVTec and the corresponding normal partition for VisA and MPDD). Source normal status is therefore used, but target test labels never enter reference construction, threshold proposal, or certification. Unless otherwise specified, $k\in\{1,2,4,8\}$ is employed for the clean accuracy-storage and strict nested studies, whereas the target-only corruption tables use $k\in\{4,8\}$. The $k=1$ strict cells use the declared patch-split support calibration fallback and are interpreted as stress tests; LOIO calibration itself requires $k\geq2$.

\subsection{Evaluation Protocols}
The target-only corruption audit evaluates all test images from all 15 MVTec and 12 VisA categories at $k\in\{4,8\}$ under five support-sampling seeds, with the $k$ sets nested within each seed, and four corruptions. Support images remain clean; each corruption is applied only to evaluation views. Images are represented on $[0,1]$. Gaussian noise adds independent $\mathcal N(0,0.05^2)$ pixel-channel noise and clips the result; blur uses a box-filter radius of one pixel (the declared three-pixel kernel); brightness/contrast applies $\operatorname{clip}(1.15(x-0.5)+0.55,0,1)$; and JPEG uses a quality-60 encode/decode. Each transformed array is clipped, quantized to 8-bit, and stored as a PNG before feature extraction. The noise generator uses the support seed plus the fixed image index, while the other transformations are deterministic. Its primary rank values follow Eq.~\eqref{eq:loio}: calibration scores use LOIO fits and each test score uses the full-support fit. The fold-matched construction is reported only as a sensitivity analysis. Alarm comparisons use the predeclared tolerance $10^{-6}$ to absorb fp32 representation error at attainable grid points.

The primary empirical groups are the target-only corruption audit and the frozen strict nested CRESS audit; the category-count feasibility curves are analytic design calculations. Supporting groups cover clean ranking and storage, PCA64/PCA128 sensitivity, direct pooled-source conformal, and attainable operating levels. The strict nested protocol uses $k\in\{1,2,4,8\}$, seeds from 0 to 4, $\alpha\in\{0.05,0.10,0.20\}$, and clean, Gaussian-noise, blur, brightness/contrast, and JPEG conditions. Its shared score export retains at most 120 evaluation records per category, seed, and condition through deterministic label-stratified sampling, so dataset labels determine the retained evaluation strata. Target labels are then used only to score evaluation outcomes; they do not enter threshold construction. Downstream source construction discards every anomalous source row, so source anomaly scores enter neither the reference map nor proposal or certification losses. The remaining cells use $\rho=0.01$, PCA64, $\delta=0.05$, and realized $M=|\mathcal T|\leq20$. Raw scores are first normalized using support-only statistics; condition-specific source views are then selected or aggregated as defined in Section~\ref{sec:cress-source-construction}. For every target category and seed, the eligible source categories are deterministically split using the frozen $0.50/0.25/0.25$ reference/proposal/certification allocation. This gives 7/3/4 categories for a 14-category within-MVTec source pool, 5/3/3 for an 11-category within-VisA pool, and 7/4/4 for the 15-category MVTec transfer pool. The same category partition is reused across conditions and $k$. The protocol evaluates MVTec and VisA within-dataset, MVTec-to-VisA, and MVTec-to-MPDD under matched-condition, clean-source, condition-agnostic, and mismatched negative-control source selection.

The frozen operational gate first requires a positive category-level threshold in at least 80\% of target cells. It then checks target FAR, power, no-harm versus LOIO, and power gain below the target-only floor. The 80\% cutoff is a predeclared usability criterion, not a confidence level; because the observed rate is zero, any positive cutoff gives the same verdict. The confidence allocation in Proposition~1 applies within each fixed target, seed, source-view mode, and condition cell; it is not a simultaneous 95\% guarantee over the complete experimental grid. Failed cells are retained.

\subsection{Baselines and Ablations}
We compare the inherited PCA residual scorer with a controlled NN memory bank using the same frozen DINOv2 features and few-shot split. PCA64 and PCA128 retain 64 and 128 principal directions, respectively; PCA64 is the default reliability substrate. For the source-assisted study, direct pooled-source conformal is an uncertified baseline: it uses the complete selected source pool with $\tau=\alpha$, but performs no reference, proposal, or certification split and applies no UCB gate. The controlled memory-bank implementation is not an official reproduction of PatchCore or AnomalyDINO; it isolates scoring and storage trade-offs under a common protocol. AnomalyDINO-S (448)~\cite{damm2025anomalydino} is included through values reported under its published protocol. WinCLIP~\cite{winclip2023} is likewise included only through published values because the audit identifies no author-released implementation associated with the cited paper. SubspaceAD is treated as a prior-art guardrail rather than a controlled baseline: the frozen DINOv2 subspace residual is inherited, and the contribution begins with the reliability analysis downstream of that score.

\subsection{Metrics}
Ranking metrics are distinguished from reliability metrics. Image AUROC and AP are calculated from the raw PCA residual score $s(x)$; pixel maps are not certification targets. Reliability is evaluated operationally by empirical FAR at fixed $\alpha$, detection power, alarm precision, the empirical cumulative distribution function (CDF) of normal rank values, and the fraction of cells receiving a positive certified threshold. Storage is quantified as retained ranker state per category. Corruption views are controlled stress tests and are not treated as independent images or as a proxy for all deployment shifts.

\subsection{Reproducibility and Claim Rules}
We conduct all experiments on a single NVIDIA RTX 5090 graphics processing unit (GPU) using PyTorch 2.12.1. DINOv2 patch features are cached in 32-bit floating-point (fp32) format. Reproducibility records include the declared protocol revision, recorded worktree state, DINOv2 source-tree and weight hashes, environment specifications, dataset counts, per-image predictions, base-image identities, support manifests, corruption parameters, source partitions, candidate-threshold bounds, and Secure Hash Algorithm 256-bit (SHA-256) checksums. Automated integrity checks pass for 1,500 MVTec, 1,200 VisA, and 600 MPDD view cells and reject duplicate cells, missing support statistics, or overlap between support and test sets. Reliability methods cannot claim ranking improvement unless the raw score is altered; image-level and category-level certificates are not interchangeable.

OpenAI Codex assisted code review, test generation, and audit scripting. The authors reviewed all changes, and no AI tool generated or altered empirical data.
